\documentclass{amsart}
% For dealing with references we use the comment environment
\usepackage{verbatim}
\newenvironment{reference}{\comment}{\endcomment}
%\newenvironment{reference}{}{}
\newenvironment{slogan}{\comment}{\endcomment}
\newenvironment{history}{\comment}{\endcomment}

% For commutative diagrams we use Xy-pic
\usepackage[all]{xy}

% We use 2cell for 2-commutative diagrams.
\xyoption{2cell}
\UseAllTwocells

% We use multicol for the list of chapters between chapters
\usepackage{multicol}

% This is generally recommended for better output
\usepackage{lmodern}
\usepackage[T1]{fontenc}

% For cross-file-references

% Package for hypertext links:
\usepackage{hyperref}

% For any local file, say "hello.tex" you want to link to please

% Theorem environments.
%
\theoremstyle{plain}
\newtheorem{theorem}[subsection]{Theorem}
\newtheorem{proposition}[subsection]{Proposition}
\newtheorem{lemma}[subsection]{Lemma}

\theoremstyle{definition}
\newtheorem{definition}[subsection]{Definition}
\newtheorem{example}[subsection]{Example}
\newtheorem{exercise}[subsection]{Exercise}
\newtheorem{situation}[subsection]{Situation}

\theoremstyle{remark}
\newtheorem{remark}[subsection]{Remark}
\newtheorem{remarks}[subsection]{Remarks}

\numberwithin{equation}{subsection}

% Macros
%
\def\lim{\mathop{\mathrm{lim}}\nolimits}
\def\colim{\mathop{\mathrm{colim}}\nolimits}
\def\Spec{\mathop{\mathrm{Spec}}}
\def\Hom{\mathop{\mathrm{Hom}}\nolimits}
\def\Ext{\mathop{\mathrm{Ext}}\nolimits}
\def\SheafHom{\mathop{\mathcal{H}\!\mathit{om}}\nolimits}
\def\SheafExt{\mathop{\mathcal{E}\!\mathit{xt}}\nolimits}
\def\Sch{\mathit{Sch}}
\def\Mor{\mathop{\mathrm{Mor}}\nolimits}
\def\Ob{\mathop{\mathrm{Ob}}\nolimits}
\def\Sh{\mathop{\mathit{Sh}}\nolimits}
\def\NL{\mathop{N\!L}\nolimits}
\def\CH{\mathop{\mathrm{CH}}\nolimits}
\def\proetale{{pro\text{-}\acute{e}tale}}
\def\etale{{\acute{e}tale}}
\def\QCoh{\mathit{QCoh}}
\def\Ker{\mathop{\mathrm{Ker}}}
\def\Im{\mathop{\mathrm{Im}}}
\def\Coker{\mathop{\mathrm{Coker}}}
\def\Coim{\mathop{\mathrm{Coim}}}

% Boxtimes
%
\DeclareMathSymbol{\boxtimes}{\mathbin}{AMSa}{"02}

%
% Macros for moduli stacks/spaces
%
\def\QCohstack{\mathcal{QC}\!\mathit{oh}}
\def\Cohstack{\mathcal{C}\!\mathit{oh}}
\def\Spacesstack{\mathcal{S}\!\mathit{paces}}
\def\Quotfunctor{\mathrm{Quot}}
\def\Hilbfunctor{\mathrm{Hilb}}
\def\Curvesstack{\mathcal{C}\!\mathit{urves}}
\def\Polarizedstack{\mathcal{P}\!\mathit{olarized}}
\def\Complexesstack{\mathcal{C}\!\mathit{omplexes}}
% \Pic is the operator that assigns to X its picard group, usage \Pic(X)
% \Picardstack_{X/B} denotes the Picard stack of X over B
% \Picardfunctor_{X/B} denotes the Picard functor of X over B
\def\Pic{\mathop{\mathrm{Pic}}\nolimits}
\def\Picardstack{\mathcal{P}\!\mathit{ic}}
\def\Picardfunctor{\mathrm{Pic}}
\def\Deformationcategory{\mathcal{D}\!\mathit{ef}}

\setlength{\emergencystretch}{3em}
\begin{document}
\title[Stacks Project, Tag 07PS]{1799 --- Every Noetherian complete local ring is a G-ring: all prime-localization completion maps have geometrically regular formal fibres}
\maketitle
\noindent Copyright (C) 2005--2025 Johan de Jong. Stacks Project authors. Source: \href{https://stacks.math.columbia.edu/tag/07PS}{Tag 07PS}.

\noindent Editorial difficulty note (not author text). Difficulty H3: Grothendieck\textquoteright{}s G-ring theorem for every complete local Noetherian ring requires geometric regularity of all localized formal fibres, not merely regularity of the original completion. The retained proof reduces through finite Cohen models, detects non-pth-powers by derivations over finite-index coefficient subfields, extends/scales those derivations through completions, and inducts through finite purely inseparable field extensions to preserve generic-fibre regularity. This positive-characteristic formal-fibre structure theorem is substantially beyond ordinary completion and regular-local-ring criteria..

\noindent Editorial provenance note (not author text). History: excerpt from revision \texttt{a0444\allowbreak 6e57e\allowbreak c1fbc\allowbreak 252a8\allowbreak 71afc\allowbreak ec775\allowbreak 2fb28\allowbreak 07b14}, more-algebra.tex, lines 12355--12389. Reference presentation was adapted; original-author context and core support blocks were retained or added. The mathematical wording is unchanged. Licensed under GNU FDL 1.2 or later; no invariant sections or cover texts. See ../licenses/COPYING. Context blocks reproduce author definitions and supporting results; they are not additional counted problems.

\begin{definition}
\label{definition-G-ring}
A ring $R$ is called a {\it G-ring} if $R$ is Noetherian and for every
prime $\mathfrak p$ of $R$ the ring map
$R_\mathfrak p \to (R_\mathfrak p)^\wedge$ is regular.
\end{definition}

\begin{definition}
\label{definition-regular}
A ring map $R \to \Lambda$ is {\it regular} if it is flat and
for every prime $\mathfrak p \subset R$ the fibre ring
$$
\Lambda \otimes_R \kappa(\mathfrak p) =
\Lambda_\mathfrak p/\mathfrak p\Lambda_\mathfrak p
$$
is Noetherian and geometrically regular over $\kappa(\mathfrak p)$.
\end{definition}

\begin{definition}
\label{definition-p-basis}
Let $p$ be a prime number. Let $k \to K$ be an extension of fields
of characteristic $p$. Denote $kK^p$ the compositum of $k$ and $K^p$
in $K$.
\begin{enumerate}
\item A subset $\{x_i\} \subset K$ is called {\it p-independent
over $k$} if the elements $x^E = \prod x_i^{e_i}$ where
$0 \leq e_i < p$ are linearly independent over $kK^p$.
\item A subset $\{x_i\}$ of $K$ is called a
{\it p-basis of $K$ over $k$} if the elements
$x^E$ form a basis of $K$ over $kK^p$.
\end{enumerate}
\end{definition}

\begin{definition}
\label{algebra-definition-coefficient-ring}
Let $(R, \mathfrak m)$ be a complete local ring.
A subring $\Lambda \subset R$ is
called a {\it coefficient ring} if the following conditions hold:
\begin{enumerate}
\item $\Lambda$ is a complete local ring with maximal ideal
$\Lambda \cap \mathfrak m$,
\item the residue field of $\Lambda$ maps isomorphically to the
residue field of $R$, and
\item $\Lambda \cap \mathfrak m = p\Lambda$, where $p$ is the characteristic
of the residue field of $R$.
\end{enumerate}
\end{definition}

\begin{definition}
\label{algebra-definition-cohen-ring}
A {\it Cohen ring} is a complete discrete valuation ring with
uniformizer $p$ a prime number.
\end{definition}

\begin{definition}
\label{algebra-definition-geometrically-regular}
Let $k$ be a field. Let $R$ be a Noetherian $k$-algebra.
The $k$-algebra $R$ is called {\it geometrically regular} over $k$ if
the equivalent conditions of Lemma \href{https://stacks.math.columbia.edu/tag/0381}{lemma geometrically regular (Tag 0381)} hold.
\end{definition}

\begin{definition}
\label{algebra-definition-regular}
A Noetherian ring $R$ is said to be {\it regular}
if all the localizations $R_{\mathfrak p}$ at primes are
regular local rings.
\end{definition}

\begin{definition}
\label{algebra-definition-regular-local}
Let $(R, \mathfrak m)$ be a Noetherian local ring of dimension $d$.
\begin{enumerate}
\item A {\it system of parameters of $R$} is a sequence of elements
$x_1, \ldots, x_d \in \mathfrak m$ which generates an ideal of
definition of $R$,
\item if there exist $x_1, \ldots, x_d \in \mathfrak m$
such that $\mathfrak m = (x_1, \ldots, x_d)$ then we call
$R$ a {\it regular local ring} and $x_1, \ldots, x_d$ a {\it regular
system of parameters}.
\end{enumerate}
\end{definition}

\begin{definition}
\label{algebra-definition-separable-field-extension}
Let $K/k$ be a field extension.
\begin{enumerate}
\item We say $K$ is {\it separably generated over $k$} if there exists
a transcendence basis $\{x_i; i \in I\}$ of $K/k$ such that the extension
$K/k(x_i; i \in I)$ is a separable algebraic extension.
\item We say $K$ is {\it separable over $k$} if for every subextension
$k \subset K' \subset K$ with $K'$ finitely generated
over $k$, the extension $K'/k$ is separably generated.
\end{enumerate}
\end{definition}

\begin{definition}
\label{algebra-definition-perfect}
Let $k$ be a field. We say $k$ is {\it perfect}
if every field extension of $k$ is separable over $k$.
\end{definition}

\begin{lemma}
\label{lemma-helper-G-ring}
Let $k$ be a field of characteristic $p$.
Let $A = k[[x_1, \ldots, x_n]][y_1, \ldots, y_m]$ and denote $K$
the fraction field of $A$.
Let $\mathfrak p \subset A$ be a prime. Then
$A_\mathfrak p^\wedge \otimes_A K$ is geometrically regular over $K$.
\end{lemma}

\begin{proof}
Let $L/K$ be a finite purely inseparable field extension.
We will show by induction on $[L : K]$ that $A_\mathfrak p^\wedge \otimes L$
is regular. The base case is $L = K$: as $A$ is regular,
$A_\mathfrak p^\wedge$ is regular (Lemma \ref{lemma-completion-regular}),
hence the localization $A_\mathfrak p^\wedge \otimes K$ is regular.
Let $K \subset M \subset L$ be a subfield such that
$L$ is a degree $p$ extension of $M$ obtained by adjoining a $p$th root
of an element $f \in M$. Let $B$ be a finite $A$-subalgebra
of $M$ with fraction field $M$. Clearing denominators, we may and do assume
$f \in B$. Set $C = B[z]/(z^p -f)$ and note that $B \subset C$
is finite and that the fraction field of $C$ is $L$. Since
$A \subset B \subset C$ are finite and $L/M/K$ are purely
inseparable we see that for every element of $B$ or $C$ some power of
it lies in $A$. Hence there is a unique prime $\mathfrak r \subset B$,
resp.\ $\mathfrak q \subset C$ lying over $\mathfrak p$. Note that
$$
A_\mathfrak p^\wedge \otimes_A M = B_\mathfrak r^\wedge \otimes_B M
$$
see Algebra, Lemma \ref{algebra-lemma-completion-finite-extension}.
By induction we know that this ring is regular. In the same manner we have
$$
A_\mathfrak p^\wedge \otimes_A L =
C_\mathfrak r^\wedge \otimes_C L =
B_\mathfrak r^\wedge \otimes_B M[z]/(z^p - f)
$$
the last equality because the completion of
$C = B[z]/(z^p - f)$ equals $B_\mathfrak r^\wedge[z]/(z^p -f)$.
By Lemma \ref{lemma-find-D} we know there exists a derivation
$D : B \to B$ such that $D(f) \not = 0$. In other words, $g = D(f)$
is a unit in $M$! By Lemma \ref{lemma-derivation-extends}
$D$ extends to a derivation of $B_\mathfrak r$, $B_\mathfrak r^\wedge$
and $B_\mathfrak r^\wedge \otimes_B M$ (successively extending through a
localization, a completion, and a localization). Since it is an
extension we end up with a derivation of $B_\mathfrak r^\wedge \otimes_B M$
which maps $f$ to $g$ and $g$ is a unit of the ring
$B_\mathfrak r^\wedge \otimes_B M$.
Hence $A_\mathfrak p^\wedge \otimes_A L$ is regular by
Lemma \ref{lemma-degree-p-extension-regular} and we win.
\end{proof}


\begin{lemma}
\label{lemma-G-ring-goes-up-quasi-finite}
Let $R \to R'$ be a finite type map of Noetherian rings and let
$$
\xymatrix{
\mathfrak q' \ar[r] & \mathfrak p' \ar[r] & R' \\
\mathfrak q \ar[r] \ar@{-}[u] &
\mathfrak p \ar[r] \ar@{-}[u] & R \ar[u]
}
$$
be primes. Assume $R \to R'$ is quasi-finite at $\mathfrak p'$.
\begin{enumerate}
\item If the formal fibre $R_\mathfrak p^\wedge \otimes_R \kappa(\mathfrak q)$
is geometrically regular over $\kappa(\mathfrak q)$, then the formal fibre
$(R'_{\mathfrak p'})^\wedge \otimes_{R'} \kappa(\mathfrak q')$
is geometrically regular over $\kappa(\mathfrak q')$.
\item If the formal fibres of $R_\mathfrak p$ are geometrically regular,
then the formal fibres of $R'_{\mathfrak p'}$ are geometrically regular.
\item If $R \to R'$ is quasi-finite and $R$ is a G-ring, then $R'$ is
a G-ring.
\end{enumerate}
\end{lemma}

\begin{proof}
It is clear that (1) $\Rightarrow$ (2) $\Rightarrow$ (3).
Assume $R_\mathfrak p^\wedge \otimes_R \kappa(\mathfrak q)$
is geometrically regular over $\kappa(\mathfrak q)$.
By Algebra, Lemma \href{https://stacks.math.columbia.edu/tag/07NC}{lemma completion at quasi finite prime (Tag 07NC)}
we see that
$$
R_\mathfrak p^\wedge \otimes_R R'
=
(R'_{\mathfrak p'})^\wedge \times B
$$
for some $R_\mathfrak p^\wedge$-algebra $B$. Hence
$R'_{\mathfrak p'} \to (R'_{\mathfrak p'})^\wedge$ is a factor of
a base change of the map $R_\mathfrak p \to R_\mathfrak p^\wedge$.
It follows that $(R'_{\mathfrak p'})^\wedge \otimes_{R'} \kappa(\mathfrak q')$
is a factor of
$$
R_\mathfrak p^\wedge \otimes_R R' \otimes_{R'} \kappa(\mathfrak q') =
R_\mathfrak p^\wedge \otimes_R \kappa(\mathfrak q)
\otimes_{\kappa(\mathfrak q)} \kappa(\mathfrak q').
$$
Thus the result follows as extension of base field preserves
geometric regularity, see
Algebra, Lemma \href{https://stacks.math.columbia.edu/tag/0381}{lemma geometrically regular (Tag 0381)}.
\end{proof}


\begin{lemma}
\label{lemma-find-D}
Let $p$ be a prime number. Let $B$ be a domain with $p = 0$ in $B$.
Let $f \in B$ be an element which is not a $p$th power in the fraction
field of $B$. If $B$ is of finite type over a Noetherian complete
local ring, then there exists a derivation $D : B \to B$ such that $D(f)$
is not zero.
\end{lemma}

\begin{proof}
Let $R$ be a Noetherian complete local ring such that there exists
a finite type ring map $R \to B$. Of course we may replace $R$ by
its image in $B$, hence we may assume $R$ is a domain of characteristic
$p > 0$ (as well as Noetherian complete local). By Algebra, Lemma
\ref{algebra-lemma-complete-local-Noetherian-domain-finite-over-regular}
we can write $R$ as a finite extension of $k[[x_1, \ldots, x_n]]$ for some
field $k$ and integer $n$. Hence we may replace $R$ by $k[[x_1, \ldots, x_n]]$.
Next, we use
Algebra, Lemma \ref{algebra-lemma-Noether-normalization-over-a-domain}
to factor $R \to B$ as
$$
R \subset R[y_1, \ldots, y_d] \subset B' \subset B
$$
with $B'$ finite over $R[y_1, \ldots, y_d]$ and $B'_g \cong B_g$
for some nonzero $g \in R$. Note that $f' = g^{pN} f \in B'$ for some
large integer $N$. It is clear that $f'$ is not a $p$th power in
the fraction field of $B'$. If we can find a derivation
$D' : B' \to B'$ with $D'(f') \not = 0$, then
Lemma \ref{lemma-derivation-extends}
guarantees that $D = g^MD'$ extends to $B$ for some $M > 0$. Then
$D(f) = g^MD'(f) = g^MD'(g^{-pN}f') = g^{M - pN}D'(f')$ is nonzero.
Thus it suffices to prove the lemma in case
$B$ is a finite extension of $A = k[[x_1, \ldots, x_n]][y_1, \ldots, y_m]$.

\medskip\noindent
Assume $B$ is a finite extension of
$A = k[[x_1, \ldots, x_n]][y_1, \ldots, y_m]$.
Denote $L$ the fraction field of $B$.
Note that $\text{d}f$ is not zero in $\Omega_{L/\mathbf{F}_p}$, see
Algebra, Lemma \ref{algebra-lemma-derivative-zero-pth-power}.
We apply Lemma \ref{lemma-power-series-ring-subfields} to find a subfield
$k' \subset k$ of finite index such that with
$A' = k'[[x_1^p, \ldots, x_n^p]][y_1^p, \ldots, y_m^p]$
the element $\text{d}f$ does not map to zero in $\Omega_{L/K'}$
where $K'$ is the fraction field of $A'$.
Thus we can choose a $K'$-derivation $D' : L \to L$
with $D'(f) \not = 0$. Since $A' \subset A$ and $A \subset B$ are
finite by construction we see that $A' \subset B$ is finite.
Choose $b_1, \ldots, b_t \in B$ which generate $B$ as an $A'$-module.
Then $D'(b_i) = f_i/g_i$ for some $f_i, g_i \in B$ with $g_i \not = 0$.
Setting $D = g_1 \ldots g_t D'$ we win.
\end{proof}


\begin{lemma}
\label{lemma-derivation-extends}
Let $R$ be a ring. Let $D : R \to R$ be a derivation.
\begin{enumerate}
\item For any ideal $I \subset R$ the derivation $D$ extends
canonically to a derivation $D^\wedge : R^\wedge \to R^\wedge$
on the $I$-adic completion.
\item For any multiplicative subset $S \subset R$ the derivation
$D$ extends uniquely to the localization $S^{-1}R$ of $R$.
\end{enumerate}
If $R \subset R'$ is a finite type extension of rings such that
$R_g \cong R'_g$ for some $g \in R$ which is a nonzerodivisor in $R'$,
then $g^ND$ extends to $R'$ for some $N \geq 0$.
\end{lemma}

\begin{proof}
Proof of (1). For $n \geq 2$ we have $D(I^n) \subset I^{n - 1}$
by the Leibniz rule. Hence $D$ induces maps $D_n : R/I^n \to R/I^{n - 1}$.
Taking the limit we obtain $D^\wedge$. We omit the verification that
$D^\wedge$ is a derivation.

\medskip\noindent
Proof of (2). To extend $D$ to $S^{-1}R$ just set
$D(r/s) = D(r)/s - rD(s)/s^2$ and check the axioms.

\medskip\noindent
Proof of the final statement. Let $x_1, \ldots, x_n \in R'$ be generators
of $R'$ over $R$. Choose an $N$ such that $g^Nx_i \in R$.
Consider $g^{N + 1}D$. By (2) this extends to $R_g$. Moreover, by
the Leibniz rule and our construction of the extension above we have
$$
g^{N + 1}D(x_i) = g^{N + 1}D(g^{-N} g^Nx_i) = -Ng^Nx_iD(g) +
gD(g^Nx_i)
$$
and both terms are in $R$. This implies that
$$
g^{N + 1}D(x_1^{e_1} \ldots x_n^{e_n}) =
\sum e_i x_1^{e_1} \ldots x_i^{e_i - 1} \ldots x_n^{e_n} g^{N + 1}D(x_i)
$$
is an element of $R'$. Hence every element of $R'$ (which can be written
as a sum of monomials in the $x_i$ with coefficients in $R$) is mapped to an
element of $R'$ by $g^{N + 1}D$ and we win.
\end{proof}


\begin{lemma}
\label{lemma-degree-p-extension-regular}
Let $R$ be a regular ring. Let $f \in R$.
Assume there exists a derivation $D : R \to R$ such that $D(f)$ is a unit
of $R$. Then $R[z]/(z^n - f)$ is regular for any integer $n \geq 1$.
More generally, $R[z]/(p(z) - f)$ is regular for any $p \in \mathbf{Z}[z]$.
\end{lemma}

\begin{proof}
By Algebra, Lemma \href{https://stacks.math.columbia.edu/tag/07NF}{lemma regular goes up (Tag 07NF)} we see that
$R[z]$ is a regular ring. Apply Lemma \ref{lemma-quotient-regular}
to the extension of $D$ to $R[z]$ which maps $z$ to zero.
This works because $D$ annihilates any polynomial with
integer coefficients and sends $f$ to a unit.
\end{proof}


\begin{lemma}
\label{lemma-quotient-regular}
\begin{slogan}
The Jacobian criterion for hypersurfaces, done right.
\end{slogan}
Let $R$ be a regular ring. Let $f \in R$. Assume there exists a
derivation $D : R \to R$ such that $D(f)$ is a unit of $R/(f)$.
Then $R/(f)$ is regular.
\end{lemma}

\begin{proof}
It suffices to prove this when $R$ is a local ring with maximal ideal
$\mathfrak m$ and residue field $\kappa$. In this case it suffices
to prove that $f \not \in \mathfrak m^2$, see
Algebra, Lemma \href{https://stacks.math.columbia.edu/tag/00NQ}{lemma regular ring CM (Tag 00NQ)}.
However, if $f \in \mathfrak m^2$ then $D(f) \in \mathfrak m$
by the Leibniz rule, a contradiction.
\end{proof}


\begin{lemma}
\label{lemma-completion-regular}
Let $A$ be a Noetherian local ring.
Then $A$ is regular if and only if $A^\wedge$ is so.
\end{lemma}

\begin{proof}
If $A^\wedge$ is regular, then $A$ is regular by
Algebra, Lemma \href{https://stacks.math.columbia.edu/tag/00OF}{lemma flat under regular (Tag 00OF)}.
Assume $A$ is regular. Let $\mathfrak m$ be the maximal ideal
of $A$. Then $\dim_{\kappa(\mathfrak m)} \mathfrak m/\mathfrak m^2 =
\dim(A) = \dim(A^\wedge)$ (Lemma \href{https://stacks.math.columbia.edu/tag/07NV}{lemma completion dimension (Tag 07NV)}).
On the other hand, $\mathfrak mA^\wedge$ is the maximal ideal of
$A^\wedge$ and hence $\mathfrak m_{A^\wedge}$
is generated by at most $\dim(A^\wedge)$ elements. Thus $A^\wedge$ is regular.
(You can also use
Algebra, Lemma \href{https://stacks.math.columbia.edu/tag/031E}{lemma flat over regular with regular fibre (Tag 031E)}.)
\end{proof}


\begin{lemma}
\label{lemma-power-series-ring-subfields}
Let $k$ be a field of characteristic $p > 0$. Let $\{x_i\}_{i \in I}$
be a $p$-basis for $k$. Let $n, m \geq 0$.
Let $K$ be the fraction field of $A = k[[x_1, \ldots, x_n]][y_1, \ldots, y_m]$.
Let $J$ be a finite subset of $I$. Consider the subfield
$k/k_J/k^p$ generated by $k^p$ and $x_i$ with $i \in I \setminus J$.
The fraction fields $K_J$ of
$$
A_J = k_J[[x_1^p, \ldots, x_n^p]][y_1^p, \ldots, y_m^p]
$$
form a family of subfields of $K$ as in
Lemma \ref{lemma-intersection-subfields}. Moreover, each of the ring extensions
$A_J \subset A$ is finite.
\end{lemma}

\begin{proof}
Since $k/k_J$ is finite, the ring extension
$k_J[[x_1^p, \ldots, x_d^p]] \subset k[[x_1, \ldots, x_d]]$ is finite
by Algebra, Lemma \href{https://stacks.math.columbia.edu/tag/0394}{lemma finite after completion (Tag 0394)}.
This implies that $A_J \to A$ is finite.

\medskip\noindent
Let us check properties (1), (2), (3) of
Lemma \ref{lemma-intersection-subfields}.
Proof of (1). For $a \in A$ we see that $a^p \in A_J$.
Hence $K^p \subset K_J$. Proof of (2).
Suppose that $f/g^p \in K$, $f, g \in A$, $g \not = 0$
is contained in $K_J$ for every choice of $J$.
Fix $J$ for the moment. Since $f/g^p \in K_J$ we can write
$f/g^p = a/b^p$ with $a \in A_J$ and $b \in A$ nonzero.
Hence $b^p f \in A_J$. For any $A_J$-derivation $D : A \to A$ we see
that $0 = D(b^pf) = b^p D(f)$ hence $D(f) = 0$ as $A$ is a domain.
Taking $D = \partial_{x_i}$ and $D = \partial_{y_j}$ we conclude
that $f \in k[[x_1^p, \ldots, x_n^p]][y_1^p, \ldots, y_d^p]$.
Applying a $k_J$-derivation $\theta : k \to k$
we similarly conclude that all coefficients of $f$
are in $k_J$, i.e., $f \in A_J$. Since it is clear that
$A^p = \bigcap\nolimits_J A_J$ where $J$ ranges over all subfields
as in the lemma we conclude $f \in A^p$ as desired.
Proof of (3). This is clear because $K_{J \cup J'} \subset K_J \cap K_{J'}$.
\end{proof}


\begin{lemma}
\label{lemma-intersection-subfields}
Let $K$ be a field of characteristic $p$. Let $\{K_\alpha\}_{\alpha \in A}$
be a collection of subfields of $K$ with the following properties
\begin{enumerate}
\item $K^p \subset K_\alpha$ for all $\alpha \in A$,
\item $K^p = \bigcap_{\alpha \in A} K_\alpha$,
\item for $\alpha, \alpha' \in A$ there exists an $\alpha'' \in A$
such that $K_{\alpha''} \subset K_\alpha \cap K_{\alpha'}$.
\end{enumerate}
Then
\begin{enumerate}
\item the intersection of the kernels of the maps
$\Omega_{K/\mathbf{F}_p} \to \Omega_{K/K_\alpha}$ is zero,
\item for any finite extension $L/K$ we have
$L^p = \bigcap_{\alpha \in A} L^pK_\alpha$.
\end{enumerate}
\end{lemma}

\begin{proof}
Proof of (1).
Choose a $p$-basis $\{x_i\}$ for $K$ over $\mathbf{F}_p$.
Suppose that $\eta = \sum_{i \in I'} y_i \text{d}x_i$ maps to zero in
$\Omega_{K/K_\alpha}$ for every $\alpha \in A$. Here the index set
$I'$ is finite. By Lemma \ref{lemma-p-basis}
this means that for every $\alpha$ there exists a relation
$$
\sum\nolimits_E a_{E, \alpha} x^E = 0,\quad a_{E, \alpha} \in K_\alpha
$$
where $E$ runs over multi-indices $E = (e_i)_{i \in I'}$ with
$0 \leq e_i < p$. On the other hand, Lemma \ref{lemma-p-basis}
guarantees there is no such relation $\sum a_E x^E = 0$ with
$a_E \in K^p$. This is a contradiction by
Lemma \ref{lemma-intersection-subfields-subspace}.

\medskip\noindent
Proof of (2). Suppose that we have a tower $L/M/K$
of finite extensions of fields. Set $M_\alpha = M^p K_\alpha$
and $L_\alpha = L^p K_\alpha = L^p M_\alpha$. Then we can first prove that
$M^p = \bigcap_{\alpha \in A} M_\alpha$, and after that prove
that $L^p = \bigcap_{\alpha \in A} L_\alpha$. Hence it suffices to
prove (2) for primitive field extensions having no nontrivial subfields.
First, assume that $L = K(\theta)$ is separable over $K$. Then
$L$ is generated by $\theta^p$ over $K$, hence we may assume that
$\theta \in L^p$. In this case we see that
$$
L^p = K^p \oplus K^p\theta \oplus \ldots K^p\theta^{d - 1}
\quad\text{and}\quad
L^pK_\alpha =
K_\alpha \oplus K_\alpha \theta \oplus \ldots K_\alpha\theta^{d - 1}
$$
where $d = [L : K]$. Thus the conclusion is clear in this case. The other
case is where $L = K(\theta)$ with $\theta^p = t \in K$, $t \not \in K^p$.
In this case we have
$$
L^p = K^p \oplus K^pt \oplus \ldots K^pt^{p - 1}
\quad\text{and}\quad
L^pK_\alpha =
K_\alpha \oplus K_\alpha t \oplus \ldots K_\alpha t^{p - 1}
$$
Again the result is clear.
\end{proof}


\begin{lemma}
\label{lemma-intersection-subfields-subspace}
Let $K/k$ be a field extension. Let $\{K_\alpha\}_{\alpha \in A}$
be a collection of subfields of $K$ with the following properties
\begin{enumerate}
\item $k \subset K_\alpha$ for all $\alpha \in A$,
\item $k = \bigcap_{\alpha \in A} K_\alpha$,
\item for $\alpha, \alpha' \in A$ there exists an $\alpha'' \in A$
such that $K_{\alpha''} \subset K_\alpha \cap K_{\alpha'}$.
\end{enumerate}
Then for $n \geq 1$ and $V \subset K^{\oplus n}$ a $K$-vector space
we have $V \cap k^{\oplus n} \not = 0$ if and only if
$V \cap K_\alpha^{\oplus n} \not = 0$ for all $\alpha \in A$.
\end{lemma}

\begin{proof}
By induction on $n$. The case $n = 1$ follows from the assumptions.
Assume the result proven for subspaces of $K^{\oplus n - 1}$.
Assume that $V \subset K^{\oplus n}$ has nonzero intersection with
$K_\alpha^{\oplus n}$ for all $\alpha \in A$. If
$V \cap 0 \oplus k^{\oplus n - 1}$ is nonzero then we win. Hence we may
assume this is not the case. By induction hypothesis we can find
an $\alpha$ such that $V \cap 0 \oplus K_\alpha^{\oplus n - 1}$
is zero. Let $v = (x_1, \ldots, x_n) \in V \cap K_\alpha^{\oplus n}$
be a nonzero element.
By our choice of $\alpha$ we see that $x_1$ is not zero.
Replace $v$ by $x_1^{-1}v$ so that $v = (1, x_2, \ldots, x_n)$.
Note that if $v' = (x_1', \ldots, x'_n) \in V \cap K_\alpha^{\oplus n}$, then
$v' - x_1'v = 0$ by our choice of $\alpha$. Hence we see that
$V \cap K_\alpha^{\oplus n} = K_\alpha v$. If we choose some
$\alpha'$ such that $K_{\alpha'} \subset K_\alpha$, then we
see that necessarily $v \in V \cap K_{\alpha'}^{\oplus n}$ (by the
same arguments applied to $\alpha'$). Hence
$$
x_2, \ldots, x_n \in
\bigcap\nolimits_{\alpha' \in A, K_{\alpha'} \subset K_\alpha} K_{\alpha'}
$$
which equals $k$ by (2) and (3).
\end{proof}


\begin{lemma}
\label{lemma-p-basis}
Let $K/k$ be a field extension. Assume $k$ has characteristic
$p > 0$. Let $\{x_i\}$ be a subset of $K$. The following are equivalent
\begin{enumerate}
\item the elements $\{x_i\}$ are $p$-independent over $k$, and
\item the elements $\text{d}x_i$ are $K$-linearly independent
in $\Omega_{K/k}$.
\end{enumerate}
Any $p$-independent collection can be extended to a $p$-basis of $K$ over $k$.
In particular, the field $K$ has a $p$-basis over $k$.
Moreover, the following are equivalent:
\begin{enumerate}
\item[(a)] $\{x_i\}$ is a $p$-basis of $K$ over $k$, and
\item[(b)] $\text{d}x_i$ is a basis of the $K$-vector space $\Omega_{K/k}$.
\end{enumerate}
\end{lemma}

\begin{proof}
Assume (2) and suppose that $\sum a_E x^E = 0$ is a linear relation
with $a_E \in k K^p$. Let $\theta_i : K \to K$ be a $k$-derivation such that
$\theta_i(x_j) = \delta_{ij}$ (Kronecker delta). Note that any $k$-derivation
of $K$ annihilates $kK^p$. Applying $\theta_i$ to the given relation we
obtain new relations
$$
\sum\nolimits_{E, e_i > 0}
e_i a_E x_1^{e_1}\ldots x_i^{e_i - 1} \ldots x_n^{e_n} = 0
$$
Hence if we pick $\sum a_E x^E$ as the relation with minimal
total degree $|E| = \sum e_i$ for some $a_E \not = 0$, then we
get a contradiction. Hence (1) holds.

\medskip\noindent
If $\{x_i\}$ is a $p$-basis for $K$ over $k$, then
$K \cong kK^p[X_i]/(X_i^p - x_i^p)$. Hence we see that
$\text{d}x_i$ forms a basis for $\Omega_{K/k}$ over $K$.
Thus (a) implies (b).

\medskip\noindent
Let $\{x_i\}$ be a $p$-independent subset of $K$ over $k$. An application
of Zorn's lemma shows that we can enlarge this to a maximal $p$-independent
subset of $K$ over $k$. We claim that any maximal $p$-independent subset
$\{x_i\}$ of $K$ is a $p$-basis of $K$ over $k$. The claim will imply
that (1) implies (2) and establish the existence of $p$-bases.
To prove the claim let $L$ be the subfield of $K$ generated by
$kK^p$ and the $x_i$. We have to show that $L = K$. If $x \in K$
but $x \not \in L$, then $x^p \in L$ and $L(x) \cong L[z]/(z^p - x)$.
Hence $\{x_i\} \cup \{x\}$ is $p$-independent over $k$, a contradiction.

\medskip\noindent
Finally, we have to show that (b) implies (a). By the equivalence of (1)
and (2) we see that $\{x_i\}$ is a maximal $p$-independent subset
of $K$ over $k$. Hence by the claim above it is a $p$-basis.
\end{proof}


\begin{lemma}
\label{algebra-lemma-complete-local-Noetherian-domain-finite-over-regular}
Let $(R, \mathfrak m)$ be a Noetherian complete local domain.
Then there exists a $R_0 \subset R$ with the following properties
\begin{enumerate}
\item $R_0$ is a regular complete local ring,
\item $R_0 \subset R$ is finite and induces an isomorphism on
residue fields,
\item $R_0$ is either isomorphic to $k[[X_1, \ldots, X_d]]$ where $k$
is a field or $\Lambda[[X_1, \ldots, X_d]]$ where $\Lambda$ is a Cohen ring.
\end{enumerate}
\end{lemma}

\begin{proof}
Let $\Lambda$ be a coefficient ring of $R$.
Since $R$ is a domain we see that either $\Lambda$ is a field
or $\Lambda$ is a Cohen ring.

\medskip\noindent
Case I: $\Lambda = k$ is a field. Let $d = \dim(R)$.
Choose $x_1, \ldots, x_d \in \mathfrak m$
which generate an ideal of definition $I \subset R$.
(See Section \href{https://stacks.math.columbia.edu/tag/00KD}{section dimension (Tag 00KD)}.)
By Lemma \href{https://stacks.math.columbia.edu/tag/0319}{lemma change ideal completion (Tag 0319)} we see that $R$
is $I$-adically complete as well.
Consider the map $R_0 = k[[X_1, \ldots, X_d]] \to R$
which maps $X_i$ to $x_i$.
Note that $R_0$ is complete with respect to the ideal
$I_0 = (X_1, \ldots, X_d)$,
and that $R/I_0R \cong R/IR$ is finite over $k = R_0/I_0$
(because $\dim(R/I) = 0$, see Section \href{https://stacks.math.columbia.edu/tag/00KD}{section dimension (Tag 00KD)}.)
Hence we conclude that $R_0 \to R$ is finite by
Lemma \href{https://stacks.math.columbia.edu/tag/031D}{lemma finite over complete ring (Tag 031D)}.
Since $\dim(R) = \dim(R_0)$ this implies that
$R_0 \to R$ is injective (see Lemma \href{https://stacks.math.columbia.edu/tag/00OJ}{lemma integral dim up (Tag 00OJ)}),
and the lemma is proved.

\medskip\noindent
Case II: $\Lambda$ is a Cohen ring. Let $d + 1 = \dim(R)$.
Let $p > 0$ be the characteristic of the residue field $k$.
As $R$ is a domain we see that $p$ is a nonzerodivisor in $R$.
Hence $\dim(R/pR) = d$, see Lemma \href{https://stacks.math.columbia.edu/tag/00KW}{lemma one equation (Tag 00KW)}.
Choose $x_1, \ldots, x_d \in R$
which generate an ideal of definition in $R/pR$.
Then $I = (p, x_1, \ldots, x_d)$ is an ideal of definition of $R$.
By Lemma \href{https://stacks.math.columbia.edu/tag/0319}{lemma change ideal completion (Tag 0319)} we see that $R$
is $I$-adically complete as well.
Consider the map $R_0 = \Lambda[[X_1, \ldots, X_d]] \to R$
which maps $X_i$ to $x_i$.
Note that $R_0$ is complete with respect to the ideal
$I_0 = (p, X_1, \ldots, X_d)$,
and that $R/I_0R \cong R/IR$ is finite over $k = R_0/I_0$
(because $\dim(R/I) = 0$, see Section \href{https://stacks.math.columbia.edu/tag/00KD}{section dimension (Tag 00KD)}.)
Hence we conclude that $R_0 \to R$ is finite by
Lemma \href{https://stacks.math.columbia.edu/tag/031D}{lemma finite over complete ring (Tag 031D)}.
Since $\dim(R) = \dim(R_0)$ this implies that
$R_0 \to R$ is injective (see Lemma \href{https://stacks.math.columbia.edu/tag/00OJ}{lemma integral dim up (Tag 00OJ)}),
and the lemma is proved.
\end{proof}


\begin{lemma}
\label{algebra-lemma-Noether-normalization-over-a-domain}
Let $R \to S$ be an injective finite type ring map. Assume $R$ is a domain.
Then there exists an integer $d$ and a factorization
$$
R \to R[y_1, \ldots, y_d] \to S' \to S
$$
by injective maps such that $S'$ is finite over $R[y_1, \ldots, y_d]$
and such that $S'_f \cong S_f$ for some nonzero $f \in R$.
\end{lemma}

\begin{proof}
Pick $x_1, \ldots, x_n \in S$ which generate $S$ over $R$.
Let $K$ be the fraction field of $R$ and $S_K = S \otimes_R K$.
By Lemma \href{https://stacks.math.columbia.edu/tag/00OY}{lemma Noether normalization (Tag 00OY)}
we can find $y_1, \ldots, y_d \in S$ such that $K[y_1, \ldots, y_d] \to S_K$
is a finite injective map. Note that $y_i \in S$ because we may pick the
$y_j$ in the $\mathbf{Z}$-algebra generated by $x_1, \ldots, x_n$.
As a finite ring map is integral (see
Lemma \href{https://stacks.math.columbia.edu/tag/00GK}{lemma finite is integral (Tag 00GK)})
we can find monic $P_i \in K[y_1, \ldots, y_d][T]$ such that
$P_i(x_i) = 0$ in $S_K$. Let $f \in R$ be a nonzero element such that
$fP_i \in R[y_1, \ldots, y_d][T]$ for all $i$. Then $fP_i(x_i)$
maps to zero in $S_K$. Hence after replacing $f$ by another
nonzero element of $R$ we may also assume $fP_i(x_i)$ is zero in $S$.
Set $x_i' = fx_i$ and let $S' \subset S$ be the $R$-subalgebra generated by
$y_1, \ldots, y_d$ and $x'_1, \ldots, x'_n$. Note that $x'_i$ is integral over
$R[y_1, \ldots, y_d]$ as we have $Q_i(x_i') = 0$ where
$Q_i = f^{\deg_T(P_i)}P_i(T/f)$ which is a monic polynomial in
$T$ with coefficients in $R[y_1, \ldots, y_d]$ by our choice of $f$.
Hence $R[y_1, \ldots, y_d] \subset S'$ is finite by
Lemma \href{https://stacks.math.columbia.edu/tag/02JJ}{lemma characterize finite in terms of integral (Tag 02JJ)}.
Since $S' \subset S$ we have $S'_f \subset S_f$ (localization is exact).
On the other hand, the elements $x_i = x'_i/f$ in $S'_f$ generate $S_f$
over $R_f$ and hence $S'_f \to S_f$ is surjective. Whence
$S'_f \cong S_f$ and we win.
\end{proof}


\begin{lemma}
\label{algebra-lemma-completion-finite-extension}
Let $R$ be a Noetherian ring. Let $R \to S$ be a finite ring map.
Let $\mathfrak p \subset R$ be a prime and let
$\mathfrak q_1, \ldots, \mathfrak q_m$ be the primes of $S$
lying over $\mathfrak p$
(Lemma \href{https://stacks.math.columbia.edu/tag/05DR}{lemma finite finite fibres (Tag 05DR)}).
Then
$$
R_\mathfrak p^\wedge \otimes_R S =
(S_\mathfrak p)^\wedge =
S_{\mathfrak q_1}^\wedge \times \ldots \times S_{\mathfrak q_m}^\wedge
$$
where the $(S_\mathfrak p)^\wedge$ is the completion with respect to
$\mathfrak p$ and the local rings $R_\mathfrak p$ and
$S_{\mathfrak q_i}$ are completed with respect to their maximal ideals.
\end{lemma}

\begin{proof}
We may replace $R$ by the localization $R_\mathfrak p$ and
$S$ by $S_\mathfrak p = S \otimes_R R_\mathfrak p$.
Hence we may assume that $R$ is a local Noetherian ring and
that $\mathfrak p = \mathfrak m$ is its maximal ideal.
The $\mathfrak q_iS_{\mathfrak q_i}$-adic completion
$S_{\mathfrak q_i}^\wedge$ is equal to the $\mathfrak m$-adic
completion by Lemma \href{https://stacks.math.columbia.edu/tag/0394}{lemma finite after completion (Tag 0394)}.
For every $n \geq 1$ prime ideals of $S/\mathfrak m^nS$ are in 1-to-1
correspondence with the maximal ideals
$\mathfrak q_1, \ldots, \mathfrak q_m$ of $S$
(by going up for $S$ over $R$, see
Lemma \href{https://stacks.math.columbia.edu/tag/00GU}{lemma integral going up (Tag 00GU)}).
Hence
$S/\mathfrak m^nS = \prod S_{\mathfrak q_i}/\mathfrak m^nS_{\mathfrak q_i}$
by Lemma \href{https://stacks.math.columbia.edu/tag/00JB}{lemma artinian finite length (Tag 00JB)}
(using for example
Proposition \href{https://stacks.math.columbia.edu/tag/00KJ}{proposition dimension zero ring (Tag 00KJ)}
to see that $S/\mathfrak m^nS$ is Artinian).
Hence the $\mathfrak m$-adic completion $S^\wedge$ of $S$ is equal to
$\prod S_{\mathfrak q_i}^\wedge$. Finally, we have
$R^\wedge \otimes_R S = S^\wedge$ by Lemma \href{https://stacks.math.columbia.edu/tag/00MA}{lemma completion tensor (Tag 00MA)}.
\end{proof}


\begin{lemma}
\label{algebra-lemma-derivative-zero-pth-power}
Let $k$ be a perfect field of characteristic $p > 0$.
Let $K/k$ be an extension.
Let $a \in K$. Then $\text{d}a = 0$ in $\Omega_{K/k}$
if and only if $a$ is a $p$th power.
\end{lemma}

\begin{proof}
By Lemma \href{https://stacks.math.columbia.edu/tag/031G}{lemma colimit differentials (Tag 031G)} we see that there exists a subfield
$k \subset L \subset K$ such that $L/k$
is a finitely generated field extension and such that
$\text{d}a$ is zero in $\Omega_{L/k}$.
Hence we may assume that $K$ is a finitely generated field extension
of $k$.

\medskip\noindent
Choose a transcendence basis $x_1, \ldots, x_r \in K$
such that $K$ is finite separable over $k(x_1, \ldots, x_r)$.
This is possible by the definitions, see
Definitions \ref{algebra-definition-perfect} and
\ref{algebra-definition-separable-field-extension}.
We remark that the result holds for the purely transcendental
subfield $k(x_1, \ldots, x_r) \subset K$.
Namely,
$$
\Omega_{k(x_1, \ldots, x_r)/k} =
\bigoplus\nolimits_{i = 1}^r k(x_1, \ldots, x_r) \text{d}x_i
$$
and any rational function all of whose partial derivatives are zero
is a $p$th power. Moreover, we also have
$$
\Omega_{K/k} =
\bigoplus\nolimits_{i = 1}^r K\text{d}x_i
$$
since $k(x_1, \ldots, x_r) \subset K$ is finite separable
(computation omitted). Suppose $a \in K$ is an element such that
$\text{d}a = 0$ in the module of differentials. By our choice of $x_i$ we
see that the minimal polynomial $P(T) \in k(x_1, \ldots, x_r)[T]$
of $a$ is separable. Write
$$
P(T) = T^d + \sum\nolimits_{i = 1}^d a_i T^{d - i}
$$
and hence
$$
0 = \text{d}P(a) = \sum\nolimits_{i = 1}^d a^{d - i}\text{d}a_i
$$
in $\Omega_{K/k}$. By the description of
$\Omega_{K/k}$ above and the fact that $P$ was the minimal
polynomial of $a$, we see that this implies $\text{d}a_i = 0$.
Hence $a_i = b_i^p$ for each $i$. Therefore by
Fields, Lemma \href{https://stacks.math.columbia.edu/tag/031V}{lemma pth root (Tag 031V)}
we see that $a$ is a $p$th power.
\end{proof}


\begin{proposition}
\label{proposition-Noetherian-complete-G-ring}
A Noetherian complete local ring is a G-ring.
\end{proposition}

\begin{proof}
Let $A$ be a Noetherian complete local ring. By
Lemma \href{https://stacks.math.columbia.edu/tag/07PN}{lemma check G ring easy (Tag 07PN)}
it suffices to check that $B = A/\mathfrak q$ has geometrically regular
formal fibres over the minimal prime $(0)$ of $B$. Thus we may assume
that $A$ is a domain and it suffices to check the condition for
the formal fibres over the minimal prime $(0)$ of $A$.
Let $K$ be the fraction field of $A$.

\medskip\noindent
We can choose a subring $A_0 \subset A$ which is a regular complete local
ring such that $A$ is finite over $A_0$, see Algebra, Lemma
\ref{algebra-lemma-complete-local-Noetherian-domain-finite-over-regular}.
Moreover, we may assume that $A_0$ is a power series ring over a
field or a Cohen ring. By Lemma \ref{lemma-G-ring-goes-up-quasi-finite}
we see that it suffices to prove the result for $A_0$.

\medskip\noindent
Assume that $A$ is a power series ring over a field or a Cohen ring.
Since $A$ is regular the localizations $A_\mathfrak p$ are regular
(see Algebra, Definition \ref{algebra-definition-regular} and the
discussion preceding it).
Hence the completions $A_\mathfrak p^\wedge$ are regular, see
Lemma \ref{lemma-completion-regular}.
Hence the fibre $A_{\mathfrak p}^\wedge \otimes_A K$ is, as a localization
of $A_\mathfrak p^\wedge$, also regular. Thus we are done if the
characteristic of $K$ is $0$. The positive characteristic case
is the case $A = k[[x_1, \ldots, x_d]]$ which is a special case of
Lemma \ref{lemma-helper-G-ring}.
\end{proof}
\end{document}
